\scanpage{90}
\deleted{5. Étude des cas \unsure{ex}}
\uline{NB} On retrouve, pour $\alpha\ne\pm2$ i.e. $\unsure{\zeta|\rho}\ne\pm1$, les représentations $\varphi^\zeta$ de $\Pi_\infty$, \inserted{du moins si $\zeta\in k$.} Donc \uline{les seuls cas nouveaux} (sur $k$ alg. clos) sont précisément les cas $\alpha=\pm2$ (i.e. $\zeta=\pm1$) qu'on avait été amené précédemment à exclure. Nous allons les étudier séparément au § suivant.
Notons dès à présent que, même si on décidait (à tort !) de ne pas s'intéresser à ces deux cas, considérés comme trop ``exotiques'', les développements en coordonnées canoniques restent intéressants, car a) ils sont valables pour tt corps, et $\unsure{=|\overline{m}}$ tout anneau, de base $k$, et \deleted{b) ils} \inserted{pas \unsure{seult}} pour $k$ alg. clos, et b) ils permettent de suivre la variation de \deleted{$\unsure{\Pi|\varphi}$} $\varphi(\alpha)$ en termes du paramètre $\unsure{\alpha,|\alpha_s}$ intrinsèque \deleted{\unsure{a}}\unsure{.|-}
% archive sheet number: 89
